\documentclass[10pt,a4paper]{amsart}
\usepackage[margin=19mm]{geometry}
\usepackage{amsmath,amssymb,mathtools}
\usepackage[scr=rsfs,cal=euler]{mathalfa}
\usepackage{xfrac}
\usepackage[unicode,hidelinks,bookmarks=false]{hyperref}
\newcommand{\ie}{i.e.\ }
\newcommand{\cf}{cf.\ }
\newcommand{\resp}{respectively\ }
\newcommand{\Subsection}{\subsection}
\providecommand{\nobreakdash}{\nobreak\hbox{-}}
\setlength{\emergencystretch}{2em}
\multlinegap0pt
\usepackage{amssymb}
\newtheorem{theorem}{Theorem}
\newtheorem{lemma}{Lemma}
\newcommand{\LHS}{\mathrm{LHS}}
\newcommand{\RHS}{\mathrm{RHS}}
\title{Nonexistence of Whirling-Knight Tours at Half Coil Count\\ for $n\equiv 4,6\pmod 8$}
\author{Shisheng Li}
\date{Source: arXiv:2605.04603v2 (CC BY 4.0)}
\begin{document}
\maketitle

\noindent\emph{Source and licence.} Nonexistence of Whirling-Knight Tours at Half Coil Count\\ for $n\equiv 4,6\pmod 8$. Version arXiv:2605.04603v2 (CC BY 4.0), distributed by arXiv under the Creative Commons Attribution 4.0 International licence, http://creativecommons.org/licenses/by/4.0/, as recorded in the arXiv licence metadata for this exact version and shown on the abstract page https://arxiv.org/abs/2605.04603v2. Full licence notice: \texttt{licenses/arxiv-2605.04603-CC-BY-4.0.txt}.\par\smallskip

\noindent\textit{Editorial scope.} Counted unit: the article's theorem thm:T2 (source0.tex lines 85--88). The article's complete original proof of it is delivered with it (source0.tex lines 463--568, one \texttt{proof} environment of 106 lines, which the article places away from the statement). No mathematics is written, repaired or re-derived: the statement and the proof are the article's own lines, in the article's own notation, and no retained line is substituted or re-flowed.  Retained uncounted as the source-local context the proof needs: lines 198--203; lines 228--231; lines 275--287.  Nothing was omitted from any retained span, so the package carries no omission record.  No retained line contains a citation, so the wrapper carries no bibliography and no bibliography entry.  No retained line contains a figure environment, an image-inclusion command or drawing code, so no image is delivered and the package declares no asset dependency.  Editorial formatting only, in addition: an AMS \texttt{amsart} preamble that restates the article's own declarations and macros used by the retained lines, namely the environments \texttt{theorem}, \texttt{lemma} on separate counters, together with the standard packages those lines need.  The retained environments print in this document as Theorem 1, Lemma 1 in order of appearance, whereas the article prints the same items with the numbering of its own full document.  The complete native source of the article as distributed in the arXiv e-print for this exact version is bundled unchanged as \texttt{sources/199873/source0.tex}.
\begin{lemma}[Reduction Lemma]\label{lem:red}
If a whirling tour with $c$~coils exists at $(n,c)$, then
$(\mathrm{LP}_{n,c})$ is feasible.  Equivalently,
$(\mathrm{LP}_{n,c})$ infeasible implies that no whirling tour
with $c$~coils on the $n\times n$ board exists.
\end{lemma}

\begin{align}
  \LHS_e &\le 0, \qquad \forall e\in E, \label{eq:LHS}\\
  \RHS &> 0. \label{eq:RHS}
\end{align}

\begin{lemma}[Geometric Lemma]\label{lem:G}
Let $n$ be even, $p=\bigl(\tfrac{n-1}{2},\tfrac{n-1}{2}\bigr)$ the
pivot, $h=n/2$, and let $w=(i,j)$ be a cell with $i\le h-2$.
\begin{itemize}
\item[(a)] If $j=h-1$ (column directly west of the pivot), then every
CCW knight arc $u\to w$ satisfies $j_u\in\{h,h+1\}$ (in particular
$j_u\ge h$, strictly east of the pivot column) and crosses the north
plumb-line.
\item[(b)] If $j=h$ (column directly east of the pivot), then every
CCW knight arc $w\to v$ satisfies $j_v\in\{h-1,h-2\}$ and crosses the
north plumb-line.
\end{itemize}
\end{lemma}

\begin{theorem}\label{thm:T2}
For every $n\equiv 4\pmod 8$ with $n\ge 4$, no $n\times n$ whirling
knight's tour with $c=n/2$ coils exists.
\end{theorem}

\begin{proof}[Proof of Theorem~\ref{thm:T2}]
Write $T_{\mathrm{even}}, T_{\mathrm{odd}}$ for the $T$-cells with
$i+j$ even or odd.  Row $i$ ($0\le i\le h-1$) contributes $T$-cells
with $j\in[h,\,n-1-i]$, a contiguous strip of length $L_i=h-i$.
Since $h=4m+2$ is even, $i+j$ at $j=h$ has the parity of $i$, and
parities alternate as $j$ increases.  For $i$ even, $L_i$ is even,
so the strip splits evenly:  $L_i/2$ cells of each parity.  For $i$
odd, $L_i$ is odd and the strip starts (at $j=h$) with parity
\emph{odd}, giving $(L_i+1)/2$ odd cells and $(L_i-1)/2$ even cells:
one extra odd cell.  The number of odd $i$ in $[0,h-1]=[0,4m+1]$ is
exactly $2m+1$, so
\begin{equation}\label{eq:parity}
  |T_{\mathrm{odd}}| - |T_{\mathrm{even}}| = 2m+1.
\end{equation}
Hence
\[
  \RHS = \sum_v\alpha_v + \sum_v\beta_v + c\gamma
       = (|T_{\mathrm{odd}}|-|T_{\mathrm{even}}|) + 2|R| - c
       = (2m+1)+2(m+1)-(4m+2) = 1>0,
\]
verifying~\eqref{eq:RHS}.

For~\eqref{eq:LHS}, recall
$\LHS_e = \alpha_v + \beta_u - \mathbf{1}_{e\text{ crosses }N}$.
Check $\LHS_e\le 0$ on every arc by case analysis (the cases below
form a covering, not a partition: arcs with $u\in R\times\{h\}\subset
T_{\mathrm{even}}$ fall under both (ii) and (v), and (v) contains the
correct argument):
\begin{itemize}
\item[(i)] $u,v\in T$:  both endpoints lie in column $j\ge h$, so
the arc stays east of the pivot column and does not cross~$N$.
The block terms cannot contribute:  $R\times\{h-1\}$ has $j<h$
so $v\notin T$ rules out $\alpha_v=+1$;  for $u\in R\times\{h\}$,
Lemma~\ref{lem:G}(b) forces $j_v\in\{h-1,h-2\}<h$, contradicting
$v\in T$.  Knight moves flip $(i+j)\bmod 2$, so the $T$-parts of
$\alpha,\beta$ give $(\alpha_v,\beta_u)\in\{(-1,+1),(0,0)\}$ and
$\LHS_e=0$.

\item[(ii)] $u\in T\setminus(R\times\{h\})$, $v\notin T$,
$v\notin R\times\{h-1\}$ (the case $u\in R\times\{h\}$ is treated
in~(v)):  $\alpha_v=0$.  If $u\in T_{\mathrm{even}}$ then
$\beta_u=0$, so $\LHS_e\le 0$.  Suppose $u\in T_{\mathrm{odd}}$
($\beta_u=+1$);
write $u=(h-1-\alpha,h+\beta)$ with $0\le\beta\le\alpha\le h-1$
(from $u\in T$) and $\alpha+\beta$ even (from $T_{\mathrm{odd}}$).
We claim the arc crosses~$N$.

\emph{Sub-case $i_u\le h-2$ (i.e.\ $\alpha\ge 1$).}
The CCW condition $\Delta_j(2\alpha+1)+\Delta_i(2\beta+1)<0$ is
satisfied by $\Delta\in\{(1,-2),(-1,-2),(-2,-1)\}$ (always for
$u\in T$), $(2,-1)$ (when $\alpha\ge 2\beta+1$), and $(-2,1)$
(when $\alpha\le 2\beta$);  the other three steps are never CCW
for $u\in T$.  For $\beta\ge 2$ each of these CCW $\Delta$'s sends
$u$ to a $v$ with $j_v\ge h$, $i_v\le h-1$, and $i_v+j_v\le n-1$
(when $v$ is on the board), so $v\in T$, contradicting case (ii).
The remaining sub-cases are:
\begin{itemize}
\item $\beta=0$ (i.e.\ $j_u=h$):  Lemma~\ref{lem:G}(b) applies, so
the arc crosses~$N$.
\item $\beta=1$ (i.e.\ $j_u=h+1$):  the only CCW arcs with $j_v<h$
are $\Delta\in\{(1,-2),(-1,-2)\}$, both with $j_v=h-1$;  the
crossing of $j=p_j$ at parameter $t^*=3/4$ has height
$i_u\pm 3/4\le h-5/4<p_i$, so the arc crosses~$N$.
\end{itemize}

\emph{Sub-case $i_u=h-1$ (i.e.\ $\alpha=0$).}  Parity then forces
$\beta=0$, so $u=(h-1,h)$.  Of its four CCW knight arcs, the one
with $v\in T$, namely $(h-1,h)\to(h-3,h+1)$, is excluded by the
case~(ii) hypothesis;  the remaining ones lying on the board are
$(h-1,h)\to(h,h-2)$ and $(h-1,h)\to(h-2,h-2)$ (always), plus
$(h-1,h)\to(h-3,h-1)$ for $h\ge 3$.  Their crossing heights at
$j=p_j$ are $h-\tfrac34$, $h-\tfrac54$, $h-2$ respectively, each
strictly below $p_i=h-\tfrac12$, so all three cross~$N$.

In every sub-case $\LHS_e=0+1-1=0$.

\item[(iii)] $v\in T$, $u\notin T$, $u\notin R\times\{h\}$:
$\beta_u=0$, $\alpha_v\in\{-1,0\}$,
$\LHS_e=\alpha_v-\mathbf{1}_{e\text{ crosses }N}\le 0$.

\item[(iv)] $v=(r,h-1)$, $r\in R$ (block head):  $\alpha_v=+1$.
$v$ sits in column $h-1$ at row $r\le 4m=h-2$, so
Lemma~\ref{lem:G}(a) applies:  the arc crosses~$N$.  We claim
$\beta_u=0$:
(a) $u\notin T_{\mathrm{odd}}$: since $i_v+j_v=r+h-1$ is odd
($r,h$ are both even) and knight moves flip parity, $u$ has even
parity;
(b) $u\notin R\times\{h\}$: such a $u$ would need $\Delta j=-1$,
hence $\Delta i=\pm 2$, but the rows of $R$ are spaced by multiples
of $4$, so $i_v-i_u=4(\ell-k)$ is never $\pm 2$.
Hence $\LHS_e=1+0-1=0$.

\item[(v)] $u=(r,h)$, $r\in R$ (block tail):  $\beta_u=+1$.
$u$ sits in column $h$ at row $r\le h-2$, so Lemma~\ref{lem:G}(b)
applies, giving $j_v\in\{h-1,h-2\}<h$ and the arc crosses~$N$.
Since $j_v<h$, $v\notin T$, ruling out $\alpha_v=-1$ from
$T_{\mathrm{even}}$;  and the $\Delta i\notin\{\pm 2\}$ argument
rules out $v\in R\times\{h-1\}$.  Hence $\alpha_v=0$ and
$\LHS_e=0+1-1=0$.

\item[(vi)] $u,v$ both outside $T$ and the blocks:
$\alpha_v=\beta_u=0$, $\LHS_e=-\mathbf{1}_{e\text{ crosses }N}\le 0$.
\end{itemize}
Combined with $\RHS=1>0$, the certificate witnesses infeasibility of
$(\mathrm{LP}_{n,n/2})$, and Lemma~\ref{lem:red} concludes the proof.
\end{proof}

\end{document}
